% MANUSCRIPT DRAFT: pending complete independent consolidated review.
% This file contains no external proof dependencies on project-private units.
\section{Finite-state estimates and upper bounds}
\label{sec:positive-proofs}

Throughout, a departure count excludes the terminal state.  For a chain
$(Y_i)_{i=0}^M$, write
\[
 J_u=\sum_{i=0}^{M-1}\mathbf1\{Y_i=u\},\qquad J=J_{Y_M}.
\]
For the binary order-two experiment, $Y_i=(X_i,X_{i+1})$,
$X_0,\ldots,X_{M+1}$ contains $M$ departures, and
$q_u=\theta_u(d(u))$ with $d(a,b)=1-a$.  Let $C_u$ count the
corresponding designated successors and put $C=C_{Y_M}$ and $q=q_{Y_M}$.
Primed counts are computed on the corrupted word.  Every bound below permits
an arbitrary initial law.  The all-state certificate, when invoked, is
$P^2(u,v)\ge\beta\Pi(v)$ for every pair of states, with $\beta>0$.
The final pair is protected, so clean and corrupted counts always select the
same actual query row.

\subsection{Universal occupation and clean prediction}

\begin{lemma}[Actual-terminal occupation]\label{lem:occupation}
For any homogeneous chain on $d$ states, any initial law, any $M\ge0$, and
any integer $w\ge1$,
\[
 \Pr(J_{Y_M}<w)\le\min\left\{1,\frac{d^2w}{M+1}\right\}.
\]
No irreducibility or stationary-start assumption is needed.
\end{lemma}
\begin{proof}
For each $t$ set $J(t)=\sum_{i<t}\mathbf1\{Y_i=Y_t\}$,
$b_t(v)=\Pr_v(J(t)<w)$, and $B_t=\max_v b_t(v)$.  If $t\ge\ell$,
the full count dominates the count in the final $\ell$-step suffix.
Conditioning on the suffix's initial state and using homogeneity gives
$B_t\le B_\ell$.  On every path, each state can occur at most $w$ times
with fewer than $w$ earlier occurrences.  Hence
$\sum_{t=0}^M\mathbf1\{J(t)<w\}\le dw$.  Consequently
\[
 (M+1)B_M\le\sum_{t=0}^M B_t
 \le\sum_v\sum_{t=0}^M b_t(v)\le d^2w.
\]
Average over the initial law and use the probability cap.
\end{proof}

\begin{lemma}[Binary clean empirical comparison]\label{lem:empirical}
For every homogeneous binary order-two source, every initial law, and
$M\ge64$,
\[
 \E\left[\mathbf1\{J>0\}\left|\frac CJ-q\right|\right]
 \le\frac{512}{\sqrt{M+1}}.
\]
The expression inside the indicator is evaluated only when $J>0$.
\end{lemma}
\begin{proof}
For each of the four rows $u$, generate an independent stack of $2M$
Bernoulli$(q_u)$ variables $\xi_{u,j}$, independently of the initial pair.
At the $j$th departure from $u$, use $\xi_{u,j}$ to choose whether the
successor is $d(u)$.  For every fixed finite path, the used stack prefixes
have exactly the product of its transition probabilities; thus this realizes
the original law, including boundary probabilities.  Write
$S_{u,j}=\sum_{i=1}^j(\xi_{u,i}-q_u)$.  Pathwise,
$C-Jq=S_{Y_M,J}$ on $J>0$.  The stacks are not asserted independent after
conditioning on $J$, $Y_M$, or any selection event.

For a centered variable in an interval of length one, its log moment
generating function has second derivative equal to a tilted variance,
which is at most $1/4$.  Integrating twice gives
$\E e^{\lambda V}\le e^{\lambda^2/8}$.  Thus
$e^{\lambda S_{u,j}-j\lambda^2/8}$ is a nonnegative supermartingale.
Stopping at its first level crossing before the deterministic time $2v$
gives, for $x>0$,
\[
 \Pr\left(\max_{j\le2v}S_{u,j}\ge x\right)
 \le e^{-\lambda x+2v\lambda^2/8}
 \le e^{-x^2/v},\qquad \lambda=2x/v.
\]
The stopping inequality here is finite: the expectation of the stopped
process is at most its initial value by summing conditional expected
increments through time $2v$.  Apply the same argument to $-S$ and take
a union over rows.  For $Z_v=\max_u\max_{j\le2v}|S_{u,j}|$,
\begin{equation}\label{eq:p-stack-tail}
 \Pr(Z_v\ge x)\le8e^{-x^2/v}.
\end{equation}
For dyadic integers $v\le M$, let $E_v=\{v\le J<2v\}$ and
$p_v=\min\{1,32v/(M+1)\}$.  Lemma~\ref{lem:occupation} gives
$\Pr(E_v)\le p_v$.  Without any independence assumption, tail integration
and $a=\sqrt{v\log(8/p_v)}$ yield
\[
 \E[Z_v\mathbf1_{E_v}]
 \le\int_0^\infty\min\{p_v,8e^{-x^2/v}\}\,dx
 \le p_va+\frac{p_vv}{2a}
 \le2p_v\sqrt{v\log(8/p_v)}.
\]
For the middle inequality, bound the Gaussian tail above $a$ by inserting
$x/a\ge1$ and integrating.  This layer contributes at most
$2p_v\sqrt{\log(8/p_v)}/\sqrt v$ to the desired risk.
Put $N=M+1$.  Summing layers with $v\ge N/32$ costs at most
\[
 \frac{2\sqrt{32\log8}}{1-2^{-1/2}}N^{-1/2}<56N^{-1/2}.
\]
If $v_0$ is the largest dyadic integer below $N/32$, then
$N/64\le v_0<N/32$ and
$\log(N/(4v_02^{-j}))\le3(j+1)$.  The remaining layers cost at most
\[
 \frac{64\sqrt3}{\sqrt{32N}}
 \sum_{j\ge0}2^{-j/2}\sqrt{j+1}
 \le\frac{64\sqrt3}{\sqrt{32N}}(1-2^{-1/2})^{-2}
 <236N^{-1/2}.
\]
Only a numerical series was extended; all stack maxima use at most $2M$
variables.  The layers partition $1\le J\le M$, including the last layer
when $M$ is dyadic.  The safe constant $512$ follows.
\end{proof}

\subsection{Return cycles and binary hitting scales}

\begin{lemma}[Finite-state return identity]\label{lem:kac}
In a finite irreducible chain with stationary law $p$, the first positive
return $T$ to $u$ has mean $\E_uT=1/p_u$.  If
$H_u=\max_v\E_vD_u$, where $D_u$ is the ordinary hitting time, zero at
$u$, then $1/p_u\le1+H_u$.
\end{lemma}
\begin{proof}
Irreducibility and finiteness give a fixed block length and a positive
minimum probability of hitting $u$ within a block, uniformly over starting
states.  Repeated conditioning gives geometric tails for hitting and return
times, so the following expectations are finite.  Let
$\nu(v)=\E_u\sum_{j=0}^{T-1}\mathbf1\{Y_j=v\}$.  First-step conditioning
inside this sum gives
\[
 (\nu P)(v)=\E_u\sum_{j=1}^{T}\mathbf1\{Y_j=v\}=\nu(v),
\]
since $Y_0=Y_T=u$.  A finite irreducible chain has a unique stationary
probability: for completeness, if $\eta$ is another one, the ratios
$\eta(v)/p_v$ are harmonic for the reversed transition matrix; a maximal
ratio propagates along every positive reversed edge and hence is constant.
Thus $\nu=(\E_uT)p$.  Exactly one visit to $u$ occurs before $T$, so
$\nu(u)=1$, proving the identity.  A positive return is one step followed
by an ordinary hit, proving $\E_uT\le1+H_u$.
\end{proof}

\begin{lemma}[Support and weighted hitting scale]\label{lem:hitting-scale}
A binary order-two source with a positive all-state certificate has either
full stationary support or a singleton absorbing diagonal support.  In the
full-support case, with $q=(a,b,e,\tau)$ and
$\overline H_u=\max\{H_u,1/\Pi(u)\}$,
\begin{equation}\label{eq:p-weighted-hitting}
 \sum_u\Pi(u)q_u\overline H_u\le\frac{40}{\sqrt\beta}.
\end{equation}
\end{lemma}
\begin{proof}
Stationarity of consecutive pair marginals gives
$\Pi=(y,x,x,w)$, and stationary flow into each diagonal gives
\begin{equation}\label{eq:p-flow}
 ya=xe,\qquad w\tau=xb.
\end{equation}
If $x>0$, the certificate entries
$P^2(10,01)=ea\ge\beta x$ and
$P^2(01,10)=b\tau\ge\beta x$ force $a,b,e,\tau>0$.
The flow identities then imply $y,w>0$.  If $x=0$, the stationary law
is on diagonals; each supported diagonal is absorbing by stationarity.
Two absorbing supported diagonals contradict the positive certificate.
This proves the classification.  Summing a certificate row gives
$\beta\le1$.  Full support and the certificate imply irreducibility.

Assume full support and put $\gamma=(a,1,1,\tau)\ge q$.
The two certificate entries just used, together with
\eqref{eq:p-flow}, also give
\begin{equation}\label{eq:p-square-entry}
 e^2\ge\beta y,\qquad b^2\ge\beta w.
\end{equation}
Here, for example, $ea\ge\beta x=\beta ya/e$ and $a,e>0$.
For target $00$, write $h_{01},h_{10},h_{11}$ for the mean hit times.
The first-step equations are
\[
 h_{01}=1+(1-b)h_{10}+bh_{11},\quad
 h_{10}=1+(1-e)h_{01},\quad h_{11}=\tau^{-1}+h_{10}.
\]
They give
\[
 h_{01}=\frac{b+2\tau}{e\tau},\quad
 h_{10}=h_{01}-\frac b\tau-1,\quad
 h_{11}=h_{01}+\frac{1-b}{\tau}-1,
\]
and therefore $H_{00}\le b/(e\tau)+2/e+1/\tau+1$.
For target $01$, the three hit times from $00,10,11$ are
$1/a$, $1+e/a$, and $1/\tau+1+e/a$, so
$H_{01}\le2/a+1/\tau+2$.
Exchanging symbols gives the other two bounds:
\[
 H_{10}\le1/a+2/\tau+2,\qquad
 H_{11}\le e/(ba)+2/b+1/a+1.
\]
The substitution exchanges $00$ with $11$, $01$ with $10$, and
$(a,b,e,\tau)$ with $(\tau,e,b,a)$; it does not assume $b=e$.
Using \eqref{eq:p-flow} and \eqref{eq:p-square-entry},
\begin{align*}
 \sqrt\beta\,yaH_{00}
 &\le\sqrt\beta\,(w+2x+ew/b+xe)
 \le\sqrt\beta\,(w+3x)+\sqrt w\le5,\\
 \sqrt\beta\,xH_{01}
 &\le2\sqrt\beta\,y/e+\sqrt\beta\,w/b+2\sqrt\beta\,x
 \le2\sqrt y+\sqrt w+2x\le5.
\end{align*}
The symbol-exchanged inequalities give the same bound for $10,11$.
Lemma~\ref{lem:kac} implies $\overline H_u\le H_u+1$; hence
$\Pi(u)\gamma_u\overline H_u\le6/\sqrt\beta\le10/\sqrt\beta$.
Sum over four states and use $q\le\gamma$.
\end{proof}

\subsection{First hits, endpoint probabilities, and joint tails}

\begin{lemma}[Target endpoint and occupation bounds]\label{lem:endpoint-tail}
For any $d$-state chain, possibly reducible, and $v\ne u$,
\begin{equation}\label{eq:p-first-hit}
 \Pr_v(D_u=t)\le\frac d{t+1},\qquad t\ge1.
\end{equation}
For an irreducible chain, fix a target with stationary mass $p>0$,
maximum ordinary hitting mean $H$, and $h=\max\{H,1/p\}$.
For every initial state $i$ and every integer $\ell\ge1$, with
$J_u(\ell)=\sum_{j=0}^{\ell-1}\mathbf1\{Y_j=u\}$,
\begin{align}
 \Pr_i(Y_\ell=u)&\le4dp(1+h/\ell),\label{eq:p-endpoint}\\
 \Pr_i(J_u(\ell)<\ell p/16)&\le2e^{-\ell/(8h)},\label{eq:p-lower-tail}\\
 \Pr_i(Y_\ell=u,J_u(\ell)>4\ell p+1)
 &\le16dp(1+h/\ell)e^{-\ell/(2h)}.\label{eq:p-joint-tail}
\end{align}
\end{lemma}
\begin{proof}
\emph{Reciprocal return coefficients and first hits.}
For any finite stochastic matrix $P$, let $T_j$ be the $j$th positive
return to $u$, allowing defective returns, and let
$F(z)=\E_u[z^{T_1};T_1<\infty]$.  Strong Markov at returns gives
$\Pr_u(T_j=s)=[z^s]F(z)^j$.  All strong-Markov uses here can be obtained
by conditioning on each possible stopping time and state and multiplying
the subsequent transition probabilities; no conditional independent-sample
claim is involved.  Let $P_{-u}$ delete row and column $u$.  Writing $P$
in target/rest block form, direct summation over the excursion before the
first return gives
$F(z)=zP_{uu}+z^2P_{u,-u}(I-zP_{-u})^{-1}P_{-u,u}$.
The block determinant identity therefore yields
\[
 1-F(z)=\frac{\det(I-zP)}{\det(I-zP_{-u})}.
\]
These are formal series, so no convergence issue arises.  The identity
$-\log\det(I-zA)=\sum_{s\ge1}\operatorname{tr}(A^s)z^s/s$
follows by differentiating the determinant, expanding $(I-zA)^{-1}$,
and integrating from zero.  Consequently
\begin{equation}\label{eq:p-return-trace}
 \sum_{j\ge1}\frac{\Pr_u(T_j=s)}j
 =\frac{\operatorname{tr}(P^s)-\operatorname{tr}(P_{-u}^s)}s
 \le\frac ds,\qquad s\ge1.
\end{equation}
Its left side is nonnegative; the upper bound uses nonnegative entries
of $P_{-u}^s$ and diagonal entries of $P^s$ at most one.
Now change only row $u$ to send $u$ deterministically to $v\ne u$.
In this auxiliary stochastic chain, the first positive return has law
$1+D_u$ under the original start $v$.  Retain only the first-return
term of \eqref{eq:p-return-trace} to prove \eqref{eq:p-first-hit}.
The auxiliary chain need not be an order-two source.

\emph{A target-specific endpoint bound.}
Return now to the irreducible chain.  Let $g(v)=\E_vD_u$.
First-step equations and Lemma~\ref{lem:kac} give
$(Pg-g)(v)=-1+\mathbf1\{v=u\}/p$.  Telescoping expectations shows
\begin{equation}\label{eq:p-expected-count}
 \E_vJ_u(s)=p\{s+\E_vg(Y_s)-g(v)\}\le p(s+H).
\end{equation}
A loop-rotation identity is
\[
 sP^s(u,u)=\sum_v\E_v[J_u(s)\mathbf1\{Y_s=v\}]
 \le dp(s+H).
\]
Indeed, for each $j<s$, the corresponding sum is
$\sum_vP^j(v,u)P^{s-j}(u,v)=P^s(u,u)$.
For $v\ne u$, split the first-hit convolution at $t=\lfloor s/2\rfloor$.
For early hits the return factor is at most $dp(1+2H/s)$.
For late hits \eqref{eq:p-first-hit} is at most $2d/s$, and
\eqref{eq:p-expected-count} from $u$ bounds the sum of return factors
by $p(\lceil s/2\rceil+H)$.  The resulting bound is at most
$dp(2+2/s+4H/s)\le4dp(1+H/s)$.  It also holds for $v=u$;
the early sum may be empty at $s=1$.  This proves the slightly stronger
version of \eqref{eq:p-endpoint} with $H$ instead of $h$.

\emph{Return moments and positive transforms.}
The family consisting of ordinary hitting times from nontarget states and
the positive return from $u$ is closed under conditioning on survival and
the current state.  Its means are at most $h$.  For each member $V$ and
integer $r\ge1$,
\begin{equation}\label{eq:p-return-moment}
 \E V^r\le r!h^{r-1}\E V.
\end{equation}
To prove this by induction, use
$V^r\le r\sum_{j=0}^{V-1}(V-j)^{r-1}$, condition each summand on
survival through time $j$, and apply the induction hypothesis and the
uniform mean bound $h$ to the residual law.  The sum of the survival
probabilities is $\E V$.
With $\lambda=1/(4h)$, expansion in nonnegative powers gives
\begin{equation}\label{eq:p-return-transforms}
 \E_i e^{\lambda D_u}\le\frac43,\qquad
 \E_u e^{\lambda T_1}\le e^{2\lambda/p},\qquad
 f:=\E_u e^{-\lambda T_1}\le e^{-3\lambda/(4p)}.
\end{equation}
For the last inequality use $e^{-x}\le1-x+x^2/2$ and
$\E T_1^2\le2h/p$; the first two follow from the geometric-series
bound on the moments.  If $K=\lceil\ell p/16\rceil$, fewer than $K$
visits in $0,\ldots,\ell-1$ means
$D_u+T_1+\cdots+T_{K-1}\ge\ell$.
These are fresh independent return cycles after the first hit, before any
horizon restriction.  Exponential Markov and \eqref{eq:p-return-transforms}
bound this probability by
$(4/3)e^{-\lambda\ell+2\lambda(K-1)/p}
\le(4/3)e^{-7\lambda\ell/8}\le2e^{-\ell/(8h)}$.
This includes $K=1$ and proves \eqref{eq:p-lower-tail}.

\emph{A joint tail by an exact change of path measure.}
Let $h_i=\E_i e^{-\lambda D_u}$, with $h_u=1$.
The inequality $e^{-x}\ge1-x$ gives $h_i\ge3/4$ and $f\ge3/4$.
Define
\[
 Q_{ij}=e^{-\lambda}P_{ij}h_j/h_i\quad(i\ne u),
 \qquad Q_{uj}=e^{-\lambda}P_{uj}h_j/f.
\]
The hitting and return equations show that every row sums to one.
This auxiliary matrix preserves directed support, so is irreducible.
Its ordinary hitting law is the original law tilted by
$e^{-\lambda D_u}/h_i$, whence $H_Q\le2H$.
Its positive return law is tilted by $e^{-\lambda T_1}/f$, and
\[
 \frac1{2p}\le\E_QT_1\le\frac2p.
\]
For the lower bound, use
$\E[T_1e^{-\lambda T_1}]\ge1/p-\lambda(2h/p)=1/(2p)$
and $f\le1$; the upper bound uses $f\ge3/4$.
Kac for $Q$ gives $p_Q\le2p$ and
$\max(H_Q,1/p_Q)\le2h$.  The endpoint bound already proved yields
$Q^\ell(i,u)\le16dp(1+h/\ell)$.
For every length-$\ell$ path from $i$ to $u$, telescoping the factors gives
\[
 Q(\text{path})=P(\text{path})e^{-\lambda\ell}
                  f^{-J_u(\ell)}/h_i.
\]
The count is exactly the number of target departures.  Thus on
$J_u(\ell)>4\ell p+1$, summing the reverse likelihood relation gives
\[
 \Pr_i(Y_\ell=u,J_u(\ell)>4\ell p+1)
 \le h_ie^{\lambda\ell}f^{4\ell p}Q^\ell(i,u)
 \le16dp(1+h/\ell)e^{-\ell/(2h)}.
\]
This proves the joint bound with its stationary-mass factor; no product
of an unconditional upper tail and an endpoint probability was taken.
\end{proof}

\subsection{Weighted imbalance, including unsupported starts}

\begin{lemma}[Half-window imbalance]\label{lem:imbalance}
For $M\ge64$, put $n=\lfloor M/2\rfloor$ and
$J_{1,u}=\sum_{i=0}^{n-1}\mathbf1\{Y_i=u\}$.
Every binary source with a positive all-state certificate, under every
initial law, satisfies
\[
 \E\left[q_{Y_M}\mathbf1\{J>0,J_{1,Y_M}<J/256\}\right]
 \le\frac{65536}{M\sqrt\beta}.
\]
\end{lemma}
\begin{proof}
We first prove a general irreducible-chain statement.  Let $m=M-n$,
$J_{2,u}=\sum_{i=n}^{M-1}\mathbf1\{Y_i=u\}$, and assign arbitrary
weights $0\le q_u\le1$.  Write
$W_u=p_uq_u\overline H_u$, where $p_u$ is stationary mass.
For every $0<c\le1/256$ we claim
\begin{equation}\label{eq:p-generic-imbalance}
 \E\left[q_{Y_M}\mathbf1\{J_{1,Y_M}<cJ\}\right]
 \le\frac{288d}{M}\sum_uW_u.
\end{equation}
The strict imbalance event is empty at $J=0$.
Fix $u$ with $np_u\ge1$.  If
$J_{1,u}\ge np_u/16$ and $J_{2,u}\le4mp_u+1$, then $m\le2n$ gives
$J_{1,u}/(J_{1,u}+J_{2,u})\ge1/145>1/256$.
Thus imbalance requires a first-window lower tail or a second-window
upper tail.  For the first event, condition on the whole history through
$Y_n$, use the uniform future endpoint bound
\eqref{eq:p-endpoint}, and then the first-window lower tail
\eqref{eq:p-lower-tail}.  For the second event condition only on $Y_n$
and use the joint bound \eqref{eq:p-joint-tail}.  Since $m\ge n$,
\[
 \Pr_i(Y_M=u,\ J_{1,u}<cJ_u)
 \le24dp_u(1+\overline H_u/n)e^{-n/(8\overline H_u)}.
\]
After multiplication by $q_u$, this is
$(24dW_u/n)(x+1)e^{-x/8}$, where $x=n/\overline H_u$.
Since the maximum of $(x+1)e^{-x/8}$ is
$8e^{-7/8}<4$, it is at most $96dW_u/n$.
If $np_u<1$, use only the endpoint bound at time $M$.
As $W_u\ge q_u$ and $p_uq_u<q_u/n$,
\[
 q_u\Pr_i(Y_M=u)\le4d(p_uq_u+W_u/M)\le8dW_u/n.
\]
Sum over $u$, average over initial states, and use $n\ge M/3$ to prove
\eqref{eq:p-generic-imbalance}.
For full-support binary sources, Lemma~\ref{lem:hitting-scale} now gives
$288\cdot4\cdot40/(M\sqrt\beta)=46080/(M\sqrt\beta)$.

It remains to treat singleton stationary support, without discarding
unsupported actual initial states.  Suppose $00$ is absorbing, so $a=0$,
and denote the other states by $A=01$, $B=10$, $C=11$.
An original path ending in this active set has never hit $00$.
Reflect the killing edge $B\to00$ to $B\to A$, producing the auxiliary
three-state chain
\[
 R(A,B)=1-b,\quad R(A,C)=b,\quad R(B,A)=1,
 \quad R(C,B)=\tau,\quad R(C,C)=1-\tau.
\]
Every original active path has probability at most its reflected-path
probability: the only changed factor is $1-e\le1$ on $B\to A$.
The endpoint and all active-state counts are unchanged.  Thus any
nonnegative functional of these paths is dominated without conditioning
on survival.
For $b,\tau>0$ this chain is irreducible, with
\begin{equation}\label{eq:p-reflected-weights}
 Z=2\tau+b,\quad p_A=p_B=\tau/Z,\quad p_C=b/Z,
 \qquad p_Ab\overline H_A\le1,\quad
 p_C\tau\overline H_C\le1.
\end{equation}
Indeed $H_A=1+1/\tau$ and $H_C=2/b$ by first-step equations;
$\overline H_A\le2+1/\tau$ and
$\overline H_C\le1+2/b$ give the last inequalities since $b,\tau\le1$.
Apply \eqref{eq:p-generic-imbalance} with weights $(b,0,\tau)$:
the $A,C$ contribution is at most $288\cdot3\cdot2/M=1728/M$.
If $b=0$ or $\tau=0$, approximate these two parameters by positive
ones in the reflected chain.  At fixed finite horizon, path probabilities
and the weights are continuous, and the bound is uniform, so the same
conclusion follows by taking a limit.  No minorization-preserving
approximation of the original source is being asserted.
The remaining target $B$ has weight $e$, and for an active start
\begin{equation}\label{eq:p-next-absorption}
 e\Pr(Y_M=B)=\Pr(D_{00}=M+1)\le\frac4{M+2},
\end{equation}
by \eqref{eq:p-first-hit}: the sole entry edge to $00$ is $B\to00$.
Drop its imbalance event.  An initial $00$ contributes zero; mixtures
follow by linearity.  The total is at most $1732/M$, below the stated
bound because $\beta\le1$.  Exchanging symbols handles absorbing $11$
with the same fixed designated-successor rule.
\end{proof}

\subsection{Reciprocal occupation with an old visit}

\begin{lemma}[Old-visit reciprocal count]\label{lem:old-inverse}
For any homogeneous $d$-state chain and initial law, let $M\ge2$,
$0\le m\le\lfloor M/2\rfloor$, and
\[
 O_u=\sum_{i=0}^{M-m-1}\mathbf1\{Y_i=u\},\qquad O=O_{Y_M}.
\]
For every $c>0$,
\begin{align*}
 \E\frac{\mathbf1\{O\ge cJ\}}{J+1}
 &\le\frac1M\left(3d^2+2d^3
     \sum_{s=m+1}^{\lfloor(M-1)/2\rfloor}\frac1s\right)\\
 &\le\frac{3d^2+2d^3}{M}
       \left(1+\log\frac M{m+1}\right),
\end{align*}
where an empty sum is zero.  In particular, the last numerator constant
is $176$ for the pair chain.  The statement includes $J=0$ and $m=0$.
\end{lemma}
\begin{proof}
Starting at $u$, the event
$\{Y_s=u,J_u(s)=j\}$ is exactly $\{T_j=s\}$ for $j\ge1$:
the initial visit is a departure, while the terminal visit is not.
Equation~\eqref{eq:p-return-trace}, which allowed defective returns, gives
\[
 K_u(s):=\E_u\frac{\mathbf1\{Y_s=u\}}{J_u(s)+1}
 =\sum_{j\ge1}\frac{\Pr_u(T_j=s)}{j+1}\le\frac ds.
\]
Fix an initial state $v\ne u$, and put
$a_v(t)=\Pr_v(D_u=t)$.  The event $Y_M=u,J_u=0$ means $D_u=M$,
so its contribution is at most $d/(M+1)\le d/M$ by
\eqref{eq:p-first-hit}; it is impossible if $v=u$.
On $J_u>0,O_u\ge cJ_u$, integrality forces $O_u\ge1$, hence the first
hit is at most $M-m-1$.  Discarding the rest of the gate and using
strong Markov gives the exact identity
\[
 \E_v\frac{\mathbf1\{Y_M=u,O_u\ge1\}}{J_u+1}
 =\sum_{t=1}^{M-m-1}a_v(t)K_u(M-t).
\]
There is no target departure before $t$, so the full count is exactly
the post-hit departure count, with no additional unit shift.
For $t\le M/2$, use $K_u(M-t)\le2d/M$ and
$\sum_ta_v(t)\le1$; the contribution is at most $2d/M$.
This also covers $v=u$, whose expression is just $K_u(M)$.
For $t>M/2$, use $a_v(t)\le2d/M$ and write $s=M-t$;
the contribution is at most
\[
 \frac{2d^2}{M}\sum_{s=m+1}^{\lfloor(M-1)/2\rfloor}\frac1s.
\]
Together with the zero-count term this is a per-target bound.
Average over the start and sum over $d$ terminal targets to obtain the
first inequality.  Finally,
$\sum_{s=a}^b1/s\le1/a+\log(b/a)\le1+\log(M/a)$ for $b\le M$,
with the empty case automatic, gives the stated logarithmic bound.
\end{proof}

\subsection{Protected edits and the two positional gates}

\begin{lemma}[Pathwise gate reduction]\label{lem:gate}
Let $D_u$ be a count of row $u$ over a fixed set of departure indices,
and at the protected row put $D=D_{Y_M}$.  On the corrupted word use
\[
 \widehat q=\begin{cases}
 C'/J',&J'>0\ \hbox{and }512D'\ge J',\\
 0,&\hbox{otherwise}.
 \end{cases}
\]
For every allowed corruption kernel and every initial law,
\begin{align}\label{eq:p-gate-reduction}
 R\le{}&\E\left[\mathbf1\{J>0\}|C/J-q|\right]
 +48k\E\frac{\mathbf1\{D\ge J/1024\}}{J+1}
 +\E[q\mathbf1\{J>0,D<J/256\}]\notag\\
 &+\Pr(J<6156k+1)+\Pr(J<3078k+1).
\end{align}
In particular this applies to a first-half count or to the old-prefix
count reconstructed by subtracting a suffix count from the total.
\end{lemma}
\begin{proof}
One edited symbol affects at most three departure triples; each changed
triple changes the histogram's $\ell^1$ norm by at most two.  Pair-count
windows affect at most two departures.  Taking a union over the edits,
including adjacent edits and indices straddling a split, gives the safe
bounds
\[
 |J'-J|,\ |C'-C|,\ |D'-D|\le\delta:=6k.
\]
When the old count is computed by subtraction, it is still exactly a
direct count on each word, so the same bound applies without adding
two separate perturbation bounds.  Endpoint protection keeps the selected
row fixed.
Let $G'=\{J'>0,512D'\ge J'\}$,
$H=\{D\ge J/1024\}$, and
$B=\{J>0,D<J/256\}$.  On $G'\cap H$ with $J>0$,
\[
 |C'/J'-q|\le|C/J-q|+\frac{8\delta}{J+1}.
\]
For $\delta=0$ the ratios coincide.  For $\delta>0$ and $J\ge2\delta$,
$J'\ge J/2$ and ratio subtraction gives a difference at most
$4\delta/J\le8\delta/(J+1)$.  For $J<2\delta$, use the unit bound on
the difference of two probabilities.
On $G'\cap H^c$,
$D+\delta\ge D'\ge(J-\delta)/512$ whereas $D<J/1024$,
so $J<1026\delta=6156k$.
On $(G')^c\cap B^c$ with $J>0$, if $J'=0$ then $J\le\delta$;
otherwise $D'<J'/512$ and $D\ge J/256$ give
$J/256-\delta<(J+\delta)/512$, hence $J<513\delta=3078k$.
If the gate rejects on $B$, its loss is exactly $q$.
The two displayed integer-threshold tails also cover $J=0$.
These cases give \eqref{eq:p-gate-reduction} pathwise and then in
expectation, with $8\delta=48k$.  At $k=0$ the false-status cases for
$J>0$ are impossible.  Equality in the gate is accepted throughout.
Since the comparison is pathwise, arbitrary permitted attack randomization
does not change it.
\end{proof}

\begin{proposition}[Finite suffix and half-gate upper bounds]
\label{prop:finite-upper}
Let $M\ge64$, $2\le r\le M+2$, and set
\[
 m=\min\{r-2,\lfloor M/2\rfloor\},\quad
 N=M-m,\quad L=1+\log\frac M{m+1}.
\]
From the histogram obtain $J'_u,C'_u$, and from the last $m+2$ symbols
obtain $T'_u=\sum_{i=M-m}^{M-1}\mathbf1\{Y'_i=u\}$.
Use the gate in Lemma~\ref{lem:gate} with $D'_u=O'_u=J'_u-T'_u$.
This $k,\beta$-blind estimator has all-start risk
\begin{equation}\label{eq:p-suffix-upper}
 R\le\min\left\{1,\frac{65536}{\sqrt M}
       +\frac{65536}{M\sqrt\beta}
       +(8448L+147744)\frac kM\right\}.
\end{equation}
Alternatively, use $D'_u=J'_{1,u}$, the count over the first
$\lfloor M/2\rfloor$ departures.  Its all-start risk is
\begin{equation}\label{eq:p-half-upper}
 R\le\min\left\{1,\frac{65536}{\sqrt M}
       +\frac{65536}{M\sqrt\beta}+2^{23}\frac kM\right\}.
\end{equation}
The latter estimator factors exactly through the histogram, protected
pair, and the one bit $\mathbf1\{512J'_{1,Y_M}\ge J'_{Y_M}\}$.
\end{proposition}
\begin{proof}
The last $m+2$ raw symbols reveal exactly $m$ departure pairs and the
terminal pair.  Thus $O'_u$ is precisely the count over indices
$0,\ldots,N-1$, even though their raw order is not observed.
In particular $r=2$ means $m=0$ and $O'=J'$; counting $r-1$ suffix
departures would incorrectly include the terminal pair.
Since $N\ge\lceil M/2\rceil$, the clean $O$ is at least $J_1$.
Lemma~\ref{lem:imbalance} therefore bounds the bias in
\eqref{eq:p-gate-reduction} by $65536/(M\sqrt\beta)$.
Lemma~\ref{lem:old-inverse}, with $d=4,c=1/1024$, bounds its reciprocal
term by $48k\cdot176L/M=8448Lk/M$.
Lemma~\ref{lem:occupation} bounds the two tails together by
\[
 \frac{16(6156+3078)k+32}{M+1}
 =\frac{147744k+32}{M+1}.
\]
Lemma~\ref{lem:empirical} and the remaining $32/(M+1)$ are bounded by
$65536/\sqrt M$.  The unit loss cap proves \eqref{eq:p-suffix-upper}.

For the half gate we require a different reciprocal estimate, valid for
any four-state chain and start:
\begin{equation}\label{eq:p-half-inverse}
 \E\frac{\mathbf1\{J_1\ge J/1024\}}{J+1}\le\frac{130944}{M}.
\end{equation}
To prove it, fix $u$ and condition on the entire first-half history,
whose terminal state is $v$ and whose count of $u$ is $j$.
Let $m_2=M-\lfloor M/2\rfloor$ be the remaining length.  Balance implies
$J_{2,u}\le1023j$.  In the suffix, the fixed-terminal event
$Y_M=u,J_{2,u}\le1023j$ is a subset of the actual-terminal low-occupation
event.  Lemma~\ref{lem:occupation} gives conditional probability at most
$16(1023j+1)/(m_2+1)$.  Since $(J_u+1)^{-1}\le(j+1)^{-1}$,
the conditional contribution is at most $16368/(m_2+1)$, including $j=0$.
Sum over four targets and use $m_2+1\ge M/2$ to prove
\eqref{eq:p-half-inverse}.  This argument uses a deterministic prefix,
not endpoint-conditioned independent samples; Lemma~\ref{lem:occupation}
is valid at every suffix length.
Apply Lemmas~\ref{lem:gate}, \ref{lem:imbalance}, and
\ref{lem:empirical} again.  The total corruption coefficient is
\[
 48\cdot130944+16(6156+3078)=6433056<2^{23},
\]
proving \eqref{eq:p-half-upper}.
The stated bit is exactly the gate decision on every corrupted word.
If $J'=0$ its value is one, but the decoder's separate $J'>0$ test still
returns zero.  Thus neither this boundary case nor threshold equality
requires another message.  All edits occur before the exact message is
formed; the bit channel itself is not being corrupted.

At $\beta=1/M$, \eqref{eq:p-suffix-upper} gives directly, for
$g_M(k)=\max\{M^{-1/2},k/M\}$,
\[
 R\le(131072+156192L)g_M(k)
 \le287264L g_M(k)<2^{19}L g_M(k).
\]
For $r=\max\{2,\lceil\alpha M\rceil\}$ with $0<\alpha\le1$,
$m+1=\min\{r-1,\lfloor M/2\rfloor+1\}\ge\alpha M/2$.
Indeed $r-1\ge\alpha M/2$ follows by separating $\alpha M\le2$
from $\alpha M>2$.  Hence $L\le1+\log(2/\alpha)$.
Likewise \eqref{eq:p-half-upper} becomes
$\min\{1,2^{17}/\sqrt M+2^{23}k/M\}$ at the critical certificate.
\end{proof}

\subsection{A count-only oracle when the budget is supplied}

\begin{proposition}[Critical known-budget upper]\label{prop:count-oracle}
At supplied certificate $\beta=1/M$, a predictor using only the exact
triple histogram, the protected pair, and the public budget $k$ satisfies,
for every source and initial law,
\[
 R\le191488\max\{M^{-1/2},k/M\},\qquad M\ge64.
\]
For $k\ge1$ its designated probability is
\[
 \widehat q_k=\frac{(C'-6k)_+}{J'+6k+2};
\]
for $k=0$ use $(C'+1)/(J'+2)$.  This budget-dependent estimator is used
only for the oracle comparison, not for the unknown-budget results.
\end{proposition}
\begin{proof}
\emph{A weighted occupation estimate.}
For an irreducible $d$-state chain with target masses $p_u$, let
$h_u=\max\{H_u,1/p_u\}$, $0\le q_u\le1$, and $W_u=p_uq_uh_u$.
For every $A>0$ and every initial law we first prove
\begin{equation}\label{eq:p-oracle-weighted-general}
 \E\min\left\{q_{Y_M},\frac A{J+2}\right\}
 \le384d^2\frac AM+\frac{192d}{M}\sum_uW_u.
\end{equation}
Fix a target $u$.  If $h_u>M$, the endpoint bound gives
$q_u\Pr(Y_M=u)\le4dp_uq_u(1+h_u/M)\le8dW_u/M$.
If $h_u\le M$, on $J_u\ge Mp_u/48$ the minimum is at most
$48A/(Mp_u)$, and $\Pr(Y_M=u)\le8dp_u$; this costs at most
$384dA/M$.  For the complementary event, put $n=\lfloor M/2\rfloor$
and $m_2=M-n$.  Since $n\ge M/3$,
$J_u<Mp_u/48$ implies $J_{1,u}<np_u/16$.
Conditioning on the complete prefix through $Y_n$, use the uniform
endpoint bound on the remaining $m_2$ transitions, and then the
unconditional lower tail on the prefix.  The weighted contribution is
at most
\[
 8dp_uq_u(1+2h_u/M)e^{-M/(24h_u)}
 =\frac{8dW_u}{M}(x+2)e^{-x/24}
 \le\frac{192dW_u}{M},\qquad x=M/h_u.
\]
The last step uses
$\sup_{x\ge0}(x+2)e^{-x/24}=24e^{-11/12}<24$.
Summing these two cases over targets proves
\eqref{eq:p-oracle-weighted-general}.
For full-support binary sources, Lemma~\ref{lem:hitting-scale} yields
\begin{equation}\label{eq:p-oracle-weighted-binary}
 \E\min\left\{q_{Y_M},\frac A{J+2}\right\}
 \le6144\frac AM+\frac{30720}{M\sqrt\beta}.
\end{equation}

For singleton support $00$, the $00$ contribution vanishes.  Use the
reflected chain from the proof of Lemma~\ref{lem:imbalance}.  On paths
ending active, original path weights are bounded above by reflected
path weights, with exactly the same terminal state and departure counts.
Thus the nonnegative minimum functional for targets $A,C$, with weights
$(b,0,\tau)$, is dominated pathwise in the same sum over paths.
Equation~\eqref{eq:p-reflected-weights} and
\eqref{eq:p-oracle-weighted-general}, with $d=3$, bound this part by
$3456A/M+1152/M$.  The original target $B$ contributes at most
$e\Pr(Y_M=B)\le4/(M+2)$ by \eqref{eq:p-next-absorption}.
An initial $00$ contributes zero.  For boundary $b$ or $\tau$, approximate
only the reflected chain by positive parameters: there are finitely many
length-$M$ paths, and each summand is continuous in probabilities and
weights, including its minimum.  The uniform bound therefore passes to
the limit.  Symbol exchange handles singleton $11$.
Since $\beta\le1$, the larger bound
\eqref{eq:p-oracle-weighted-binary} covers both singleton cases and
arbitrary initial mixtures as well.

\emph{A pathwise lower envelope.}
For $k\ge1$ put $B_k=6k$, $D=J+2$, $D'=J'+B_k+2$, and
$c'=(C'-B_k)_+$.  The protected-edit perturbation bounds give
\[
 D'\ge D,\quad0\le c'\le C,\quad
 0\le C-c'\le2B_k,\quad0\le D'-D\le2B_k.
\]
Thus $0\le\widehat q_k\le C/(J+2)$, and
\[
 \frac CD-\frac{c'}{D'}
 =\frac{C-c'}D+\frac{c'}{D'}\frac{D'-D}D
 \le\frac{4B_k}D.
\]
For $J>0$, putting $p=C/J$ gives
\[
 0\le p-\widehat q_k
 \le\min\left\{p,\frac{4B_k+2}{J+2}\right\}.
\]
Since $x\mapsto\min\{x,c\}$ is $1$-Lipschitz on nonnegative reals,
with $A=24k+2$ the triangle inequality implies
\begin{equation}\label{eq:p-oracle-pathwise}
 |\widehat q_k-q|
 \le2\mathbf1\{J>0\}|C/J-q|
       +\min\left\{q,\frac A{J+2}\right\}.
\end{equation}
When $J=0$, $C=0$ and $\widehat q_k=0$; since $A\ge2$, the last
term equals $q$, so the inequality still holds without defining $C/J$.
This is valid for every permitted edit realization.
Using Lemma~\ref{lem:empirical},
\eqref{eq:p-oracle-weighted-binary}, $\beta=1/M$, and $A\le26k$ gives
\[
 R\le\frac{31744}{\sqrt M}+159744\frac kM
 \le191488\max\{M^{-1/2},k/M\},\qquad k\ge1.
\]

For $k=0$ the word is unedited.  On $J>0$, the add-one probability
is within $1/(J+2)$ of $C/J$; on $J=0$ pay unit loss.
Lemma~\ref{lem:occupation} gives $\Pr(J=0)\le16/(M+1)$.
For completeness, integer-threshold layer cake gives
\[
 \E\frac1{J+2}
 \le\int_0^1\min\left\{1,\frac{16}{(M+1)t}\right\}\,dt
 \le\frac8{\sqrt{M+1}}.
\]
Indeed a nonempty event $J<t^{-1}-2$ uses threshold
$\lceil t^{-1}-2\rceil\le t^{-1}$; for the last inequality use
$\min\{1,z\}\le\sqrt z$ and integrate $t^{-1/2}$.
Together with Lemma~\ref{lem:empirical} this is at most
$536/\sqrt M$, below the asserted constant.  No positional message or
raw-window input is used in either budget branch.
\end{proof}
